\frontmatterpage
\unchapter{Anotācija}

Daudzās bijušās Padomju Savienības valstīs darbojas nelieli biroji, kas maina vienu valūtu pret citu un nosūta naudu uz citām valstīm. Piemēram, ģimene Uzbekistānā var ienākt šādā birojā, lai saņemtu naudu, ko atsūtījis radinieks, kurš strādā Krievijā, un birojs reizi nedēļā norēķinās ar partneru biroju Maskavā. Šādi biroji bieži ir ģimenes uzņēmumi ar filiālēm vairākās pilsētās un valstīs. Tie ikdienā apgroza ievērojamas naudas summas, taču parasti glabā ierakstus papīra žurnālā vai Excel izklājlapā. Tāpēc kļūdas notiek bieži, un viena izlaista rinda var radīt zaudējumu vairāku tūkstošu ASV dolāru apmērā, kura izsekošanai vēlāk nepieciešamas vairākas nedēļas.

Kļūdas notiek tāpēc, ka birojam patiesībā vienlaikus ir jāuztur trīs dažādu veidu uzskaite. Pirmā ir nauda paša biroja kasē un uz paša biroja kartēm. Otrā ir tekošais parāds ar katru partneru biroju ārzemēs, un šis parāds ir jāseko vairākās valūtās vienlaikus, jo vienam un tam pašam partnerim vienlaikus var būt parāds gan Krievijas rubļos, gan Turcijas lirās. Trešā ir pastāvīgo klientu nauda, kas bieži uzkrāj birojā nelielas summas dažādās valūtās. Standarta grāmatvedības programmas visus trīs šos veidus apvieno vienā skaitlī vienā valūtā, un tas tieši slēpj to, kas birojam patiešām jāredz.

Šajā bakalaura darbā ir aprakstīta programma \emph{Ethnocount}, ko autors izstrādājis šīs reālās problēmas risināšanai. Programma strādā uz telefoniem (Android, iPhone) un datoriem (Windows, macOS, Linux) no vienota koda, tāpēc filiāle var lietot to pašu rīku neatkarīgi no tā, vai operators sēž pie galda vai pārvietojas ar telefonu rokā. Programma trīs uzskaites veidus saglabā atsevišķi --- tieši tā, kā tie patiesībā pastāv; tā pareizi summē naudu pat tad, ja vienam partnerim ir parādi trīs dažādās valūtās; un tā nodrošina, ka pārvedumu nevar nejauši saglabāt divreiz, pat ja operators sliktā interneta savienojuma dēļ nospiež pogu divas reizes. Pastāv arī noteikumi par to, kurš ko drīkst darīt: tīkla dibinātājs redz visu, direktors redz visas filiāles, bet parastais darbinieks redz tikai savu filiāli un pirms pagātnes ierakstu maiņas ir jāprasa direktora atļauja.

Autors uzbūvēja sistēmu, rūpīgi pārbaudīja to ar daudziem automātiskiem testiem, izmērīja, cik ātri tā darbojas uz parasta Windows datora un uz vidējas klases Android telefona, un izmēģināja to divas nedēļas kopā ar reālu maiņas biroju. Izmēģinājuma laikā birojs uzturēja gan veco Excel izklājlapu, gan jauno programmu, un katras nedēļas beigās tās tika salīdzinātas. Jaunā programma neradīja nevienu savu kļūdu. Gluži pretēji, tā atrada trīs vecas kļūdas izklājlapā, kuras operators iepriekš nebija pamanījis.

Darbam joprojām ir ierobežojumi, kas tiek godīgi izklāstīti. Programma vēl nepatstāvīgi neielādē valūtas kursus no interneta --- operatoram tie jāievada manuāli. Tā neļauj operatoram saglabāt pārvedumu, kad nav interneta savienojuma --- šādā gadījumā darbība tiek atteikta, jo ar naudu drošāk ir atteikt nekā riskēt ar nepareizu kārtību vēlāk. Un tā tika pārbaudīta vienā birojā divu nedēļu laikā, nevis daudzos birojos vairāku mēnešu garumā. Šie ierobežojumi norāda dabiskus nākamos soļus projekta attīstībai, kas ir uzskaitīti noslēguma nodaļā.

\thesisscopeLV

\textbf{Atslēgvārdi:} starpplatformu mobilā izstrāde, Flutter, BLoC, Supabase, PostgreSQL, Row-Level Security, glabātās procedūras, divkāršā ieraksta grāmatvedība, naudas pārvedumi, ārvalstu valūtas, vairāku filiāļu operācijas, vairāku nomnieku arhitektūra, idempotence.
